\section{Notification System Architecture}

This section describes the overall architecture of the notification system, including the message routing strategy, the interaction between system components, and the privacy considerations applied to the design.

\subsection{Overview}

The system follows a publish-subscribe architecture, where internal services produce events that are routed through RabbitMQ and eventually delivered to clients as notifications. The architecture is designed around three main principles:

\begin{itemize}
    \item \textbf{Decoupling:} Producers and consumers are independent, communicating only through the messaging system.
    \item \textbf{Flexible filtering:} Notifications are routed based on user-defined filters using content-based routing.
    \item \textbf{Data minimization:} Sensitive information is removed before events are published to the broker.
\end{itemize}

Figure~\ref{fig:architecture} illustrates the full system.

\begin{figure}[h]
    \centering
    \includegraphics[width=\textwidth]{architecture.png}
    \caption{Notification architecture with hybrid routing and data minimization}
    \label{fig:architecture}
\end{figure}

\subsection{Internal Event Production}

The system begins with internal services (e.g., a listing service) that generate events such as the creation or update of a listing. These services operate on the full internal data model, which includes sensitive information such as user or owner identity.

This internal representation must not be directly exposed to the messaging system.

\subsection{Sanitization and Event Transformation}

Before an event is published to RabbitMQ, it is processed by a dedicated sanitization layer. This component is responsible for transforming internal events into \textit{notification-safe events}.

The transformation includes:
\begin{itemize}
    \item Removing all personal or sensitive data
    \item Replacing internal references with pseudonymous identifiers (e.g., \texttt{listing\_id})
    \item Extracting only the attributes required for routing (e.g., city, country, price bucket)
\end{itemize}

This step enforces data minimization and ensures that the messaging layer does not expose private information.

\subsection{Message Routing in RabbitMQ}

After sanitization, events are published to RabbitMQ. The routing logic follows a hybrid approach combining topic-based and content-based routing.

\subsubsection{Topic-based Routing (Coarse Classification)}

Messages are first sent to a topic exchange using a routing key that represents the event type, such as:

\begin{itemize}
    \item \texttt{listing.created}
    \item \texttt{listing.updated}
    \item \texttt{listing.price\_changed}
\end{itemize}

This stage provides coarse-grained classification, separating messages based on what type of event occurred.

\subsubsection{Content-based Routing (Fine-grained Filtering)}

Messages are then forwarded to a header exchange, where routing decisions are made based on message attributes. These attributes are included as headers and may include:

\begin{itemize}
    \item \texttt{country}
    \item \texttt{city}
    \item \texttt{town}
    \item \texttt{price\_bucket}
\end{itemize}

Queues are bound to the header exchange with specific filter conditions. Only messages matching those conditions are delivered to each queue.

This mechanism allows the system to implement flexible notification filters directly at the messaging layer.

\subsubsection{Hybrid Routing Benefits}

The combination of topic-based and content-based routing enables a separation of concerns:

\begin{itemize}
    \item Topic routing answers: \textit{what type of event occurred?}
    \item Content-based routing answers: \textit{which notifications should receive it?}
\end{itemize}

Although both event types may ultimately reach the same queues in simple deployments, this structure allows future evolution toward independent routing paths, improved scalability, and better workload isolation.

\subsection{Queues and Notification Workers}

Each queue represents a set of messages matching a particular filter. Notification workers are consumer services that subscribe to these queues.

A notification worker is responsible for:

\begin{itemize}
    \item Consuming messages from a queue
    \item Optionally enriching or formatting the notification
    \item Delivering the notification to end users (e.g., via WebSockets or push services)
    \item Acknowledging successful processing
\end{itemize}

RabbitMQ itself does not deliver messages to clients directly; instead, workers act as the bridge between the messaging system and user-facing delivery mechanisms.

Multiple workers can consume from the same queue, allowing horizontal scaling of notification delivery.

\subsection{Client Delivery}

After processing, notifications are sent to client applications. Clients receive only sanitized and public-safe information, such as:

\begin{itemize}
    \item Listing identifier
    \item Location information
    \item Price range
    \item Notification message
\end{itemize}

If additional details are required, clients must retrieve them through authenticated API calls.

\subsection{Privacy and Trust Boundaries}

The architecture enforces clear separation between different trust zones:

\begin{itemize}
    \item \textbf{Internal zone:} Contains sensitive data (e.g., user identities, ownership information)
    \item \textbf{Broker zone:} Handles routing using sanitized data only
    \item \textbf{Client zone:} Receives only public or non-sensitive information
\end{itemize}

The mapping between a listing and its owner remains exclusively within internal services and is never exposed to RabbitMQ or clients.

This approach provides data minimization and pseudonymization, reducing the risk of data leakage. However, as discussed in Section~X, centralized publish-subscribe systems do not provide strong anonymity guarantees, as the broker may still observe communication patterns.

\subsection{Scalability Considerations}

The architecture supports scalability at multiple levels:

\begin{itemize}
    \item \textbf{Queue-level scaling:} Multiple workers can consume from the same queue
    \item \textbf{Routing-level scaling:} Topic-based routing allows separation of event types into different processing paths
    \item \textbf{Service-level scaling:} Notification workers can be independently scaled based on workload
\end{itemize}

The hybrid routing model enables future extensions where different event types (e.g., \texttt{created} vs \texttt{updated}) can be routed to separate exchanges or queues, allowing independent scaling and isolation of workloads.

\subsection{Summary}

The proposed architecture combines flexible message routing with strong data minimization principles. RabbitMQ is used as a routing and delivery backbone, while sensitive data remains confined to internal services. The hybrid routing approach ensures both expressiveness and scalability, making the system adaptable to evolving requirements.
